\section{Parallel Planning Approach}
\label{sec:parallel_planning_approach}
\hspace{\parindent}
The proposed architecture extends the TEA* planner by formalizing and extending the
multi-threaded planning idea previously explored in~\cite{ribeiro2025parallel}. In that initial
implementation, parallelism was used to evaluate multiple robot-solving orders in parallel, thereby 
reducing the dependence of TEA* on a single priority permutation. The framework developed in this dissertation preserves this principle, but replaces the previous
lexicographic permutation selection strategy with heuristic-guided priority orderings and integrates
the execution into a dedicated parallel planner structure.

As illustrated in Figure~\ref{fig:parallel_planning_framework}, the Supervisor sends a planning
request to the Path Planner containing the current robot positions, their corresponding goals, and
the maximum time available for planning, denoted by
\(I=\{(p_i,g_i)\}_{i=1}^{N}, T_{\max}\).

After receiving the request, the Path Planner initializes the execution and launches a set of worker
threads. Each worker retrieves a heuristic from a shared heuristic pool, uses it to compute a
priority permutation, and then runs TEA* independently using that ordering. Each execution
therefore produces a candidate plan associated with one heuristic-generated priority ordering.

\begin{figure}[H]
    \centering
    \includegraphics[width=1.0\linewidth]{parallel_planning_framework.png}
    \caption{Overview of the proposed parallel planning framework.}
    \label{fig:parallel_planning_framework}
\end{figure}

For simplicity, Figure~\ref{fig:parallel_planning_framework} represents the case in which the
number of available heuristics is equal to the number of worker threads. In that simplified case,
each thread retrieves one heuristic and produces one candidate plan. In the general case, however,
these numbers do not need to be equal, as a thread may evaluate several heuristics during the same
planning cycle while planning time is still available and heuristics remain to be tested.

During the parallel execution stage, each worker thread generates one or more candidate plans,
which are collected in a set \(P\) and subsequently evaluated using the selected cost function \(J\).
Once all candidate plans have been evaluated, the plan with the lowest cost is returned to the
Supervisor. In this work, \(J(P_i)\) is defined as the sum of the individual costs of all robot paths
in candidate plan \(P_i\). Candidate plans that do not include valid paths for all robots are assigned
an \(\infty\) cost, as a plan is considered to have failed whenever at least one robot is unable to
obtain a valid path.

The following sections describe the main components of the proposed framework in more detail,
including the worker thread implementation, the heuristic dispatching strategy, the timeout and
controlled termination mechanism, and the handling of redundant heuristics.



\section{Worker Thread Implementation Details}
\label{sec:parallel_implementation_details}
\hspace{\parindent}
The implementation of the parallel framework was based on the multi-threaded TEA* prototype
presented in~\cite{ribeiro2025parallel}, where OpenMP was used to evaluate multiple robot orderings
in parallel. In that initial implementation, however, the parallel execution logic was closely
integrated into the Path Planner code. This made the approach suitable as a proof of concept, but
less modular for the integration of the heuristic framework developed in this dissertation.

The redesigned implementation exploits the Object-Oriented Programming (OOP) structure of the
existing planner. Since TEA* was already implemented as a class with its own internal planning
state and methods, the parallel execution logic was reorganized around a dedicated TEA* Parallel
Planner class, illustrated in Figure~\ref{fig:parallel_planner_class_diagram}. This class encapsulates
the parallel coordination logic and manages vectors of independent TEA* and A* instances. The
TEA* instances are used by the worker threads to evaluate different priority orderings, while the
A* instances support heuristics that require single-robot shortest-path computations. This structure
keeps the original sequential behavior of TEA* inside each worker, while separating the parallel
execution logic from the standard Path Planner implementation.

\begin{figure}[H]
    \centering
    \includegraphics[width=1.0\linewidth]{class_diagram.png}
    \caption{UML Class Diagram representing the implemented parallel planning framework and its main components.}
    \label{fig:parallel_planner_class_diagram}
\end{figure}

Different thread-based parallelization options could be considered for this type of implementation,
including \textit{OpenMP}, \textit{Intel TBB}, and the C++ standard threading library. These
approaches rely on kernel-level threads that can be scheduled on multiple processor cores, but they
expose different programming models. \textit{Intel TBB} provides flexible task-based parallelism,
while the C++ standard threading library offers explicit thread management through
\textit{std::thread}.

Since the previous prototype already relied on \textit{OpenMP}, this choice was re-evaluated before
developing the reorganized framework. The workload considered in this dissertation remains
naturally expressed as a set of independent worker executions, each evaluating one priority
ordering at a time. Therefore, the additional control provided by alternatives such as
\textit{std::thread} or \textit{Intel TBB} was not necessary for the implemented design. For this reason, \textit{OpenMP} was adopted in the redesigned framework, as it provides a simple
and sufficiently expressive interface for launching the main worker threads while avoiding explicit
management of thread creation, joining, and scheduling.

During each planning cycle, the worker executions are independent but not equivalent. As each
worker evaluates a different priority ordering, each TEA* instance may produce a different sequence
of reservations, node expansions, and path rejections. To avoid interference between parallel
executions, the mutable internal state of TEA* is kept private to each planner instance. This includes
structures such as the 3D costmap, where the space-time reservations of the planned robot paths are
stored, as well as the lists used during the search.

This design reduces the need for synchronization during the TEA* search itself. Sharing these
structures between threads would require frequent mutex-protected accesses in some of the most
computationally intensive parts of the algorithm, increasing implementation complexity. More
importantly, workers would repeatedly block each other while accessing shared planning state,
which could largely offset the performance benefit of executing multiple TEA* instances in
parallel. However, the main drawback is the additional memory usage,
as each TEA* instance maintains its own 3D costmap, which can be large because it consists of
\(T\) roadmap representations, where \(T\) is the maximum number of planning steps considered.
Therefore, using multiple worker threads implies storing this structure multiple times, making the
memory cost proportional to the number of TEA* instances. This overhead can become significant
for large maps or large values of \(T\).

For this reason, only data that remains constant during planning is shared between instances. For
example, the precomputed map cost lookup table is shared, as it is read-only and can be safely
accessed by multiple threads without synchronization. This provides a compromise between thread
isolation, which avoids frequent synchronization during the search, and memory efficiency, by
sharing only the data that can be safely reused across workers.

Overall, the implementation combines the simplicity of \textit{OpenMP} for launching the main worker
threads with explicit synchronization mechanisms only for the few shared structures that require
controlled access. This design keeps the core TEA* algorithm unchanged, while enabling several
priority orderings to be evaluated concurrently through independent planner instances.


\section{Heuristic Dispatching Mechanism}
\label{sec:heuristic_dispatching}
\hspace{\parindent}
The previous multi-threaded TEA* prototype showed that evaluating multiple robot priority
orderings in parallel can reduce the dependence of TEA* on a single ordering
~\cite{ribeiro2025parallel}. Building on that idea, the framework developed in this dissertation uses
the prioritization heuristics introduced in Chapter~\ref{chap:ch5} to generate more informed
orderings. The dispatching mechanism was therefore designed to assign these heuristic-generated
orderings to worker threads while avoiding the simultaneous evaluation of highly redundant
heuristics.

As the heuristic set is organized into different families, the objective is to promote a balanced use
of these families during parallel execution. This increases the diversity of the generated priority
orderings and reduces the probability that several threads evaluate very similar planning
alternatives. The overall dispatching flow executed by each worker thread is illustrated in
Figure~\ref{fig:heuristic_dispatching_flow}.

To achieve this, the available heuristics are stored in a shared priority queue. Each heuristic is
associated with a heuristic group, and each group has a counter indicating how many heuristics from
that group have already been selected during the current planning cycle. The priority queue is
ordered according to these group counters, so that heuristics belonging to the least-used groups are
given higher priority.

The set of heuristics inserted into the pool is defined through the \textit{heuristics.yaml} file,
represented in Figure~\ref{fig:parallel_planner_class_diagram}. This configuration file specifies
which heuristics are available during the planning cycle and associates each one with a heuristic
group. Within each group, the relative order of the heuristics follows the order defined in the YAML
file. The dynamic component of the dispatcher is therefore applied mainly at the group level, since
heuristics belonging to the same group tend to use similar criteria to generate priority orderings.

Whenever a worker thread needs a new heuristic, it accesses the shared queue inside a critical
section protected by a mutex and pops a heuristic from the group with the lowest usage counter.
The selected heuristic is removed from the queue during this operation, preventing two threads from
evaluating the exact same heuristic during the same planning cycle. After the heuristic is selected,
the counter of the corresponding group is incremented, and the queue is reordered accordingly.

This mechanism encourages the simultaneous evaluation of heuristics from different families.
Instead of assigning several heuristics from the same family to different threads at the same time,
the dispatcher tends to distribute the available computational resources across the available
strategies.

As depicted in Figure~\ref{fig:heuristic_dispatching_flow}, if the shared heuristic queue becomes
empty, the worker can no longer retrieve a heuristic from the pool. In this case, if that worker has
not yet obtained a valid candidate plan, an additional attempt is performed using the
\textit{Exploration Conflicts Pressure} replanning heuristic. The statistics required by this heuristic
are accumulated locally by each worker across its previous planning attempts, namely the number
of reservation-based rejections and the number of expanded nodes already mentioned in
Subsection~\ref{subsec:exploration_conflicts_pressure}. This heuristic is only used after all regular
heuristic attempts of that worker have failed, as it tends to reverse the priority order in an informed
way.

Although a predefined static ordering with manually alternated heuristic families could provide a
similar behavior in the current implementation, the main advantage of the implemented mechanism
is the flexibility it provides for extending the scheduling policy. In future extensions, the dispatcher
could prioritize heuristics according to the remaining planning time or even the quality of candidate
plans obtained by previous orderings.

After leaving the critical section, the worker thread evaluates the selected heuristic by generating
the corresponding priority ordering and computing the associated TEA* paths. If the regular
heuristic queue is empty, the worker either returns the best available plan, when it has already
found a valid one, or performs an additional attempt with the \textit{Exploration Conflicts Pressure}
heuristic. The worker only repeats this procedure after checking whether the current planning cycle
is still active, according to the timeout and termination logic described in the following section.

\begin{figure}[H]
    \centering
    \includegraphics[width=1.0\linewidth]{fluxogram.drawio.png}
    \caption{Heuristic dispatching flow executed by each worker thread.}
    \label{fig:heuristic_dispatching_flow}
\end{figure}




\section{Timeout Handling and Controlled Termination}
\label{sec:timeout_controlled_termination}
\hspace{\parindent}
The parallel planner was designed to operate under a maximum planning time \(T_{\max}\), which
represents the time available to compute and return a valid plan after a planning request is received.
In the intended system behavior, this value is provided by the Supervisor and may be computed
dynamically in future versions according to the current execution state of the robots. For instance,
the Supervisor could estimate \(T_{\max}\) from the remaining time before the most urgent robot
reaches its next relevant node. On the Path Planner side, the timeout mechanism is already
supported, although in the current implementation \(T_{\max}\) is provided as a static parameter.

Immediately before launching the worker threads, a global timer is initialized. To monitor the
elapsed planning time, an additional lightweight monitoring thread is launched alongside the worker
threads. During the planning cycle, this thread periodically compares the elapsed time with
\(T_{\max}\), sleeping between consecutive checks in order to reduce computational overhead and
avoid unnecessary CPU usage. Once the elapsed time exceeds the maximum planning time, the
monitoring thread sets a shared atomic boolean flag to indicate that the timeout has been reached.

However, the timeout does not immediately force all workers to stop. Returning no plan is
undesirable when no valid solution has yet been found, so the termination condition also depends
on the existence of at least one complete candidate plan. For this reason, another shared atomic
flag is used to indicate whether a valid plan has already been obtained. The first worker thread
that finds a complete plan, containing valid paths for all robots, sets this flag.

The planning cycle is therefore considered inactive only when both of the following conditions are
satisfied:

\begin{itemize}
    \item The maximum planning time \(T_{\max}\) has been reached;
    \item At least one valid candidate plan is available.
\end{itemize}

To allow running TEA* instances to stop in a controlled manner, the termination condition is
checked at the beginning of each TEA* step. This point was chosen because the step operation is
the most frequently executed operation throughout the search process, making it a natural location
to test whether the current planning attempt should continue. When the termination condition
becomes true, each worker detects it at the beginning of a subsequent step and interrupts its
current planning process safely, allowing local variables and internal planner state to be closed
consistently before the thread returns.

This mechanism allows the planner to use the available time to continue searching for better
candidate plans after the first valid solution is found, while still ensuring that the computation can
be interrupted once the time budget has expired. As a result, the framework does not stop at the
first feasible plan, but instead uses the remaining time to improve the final selection whenever
possible.


\section{Handling Redundant Heuristics}
\label{sec:redundant_heuristics}
\hspace{\parindent}
The results presented in the previous chapter showed that \textit{Radius Neighbor Density} and
\textit{\(K\)-hop Neighbor Density} often produce equivalent or very similar priority orderings.
Running both heuristics may therefore lead to redundant TEA* executions. For this reason, the
selection between them was based on one of their main differentiating factors, the estimated
computation time. This criterion is directly motivated by their different computational behaviors,
described in Equations~\ref{eq:radius_neighbor_density_complexity} and
\ref{eq:k_hop_neighbor_density_complexity}.

For this purpose, linear regression models were fitted for both heuristics. The theoretical
complexity analysis had already identified the main runtime trends expected in each case, making
more complex modeling approaches unnecessary. The quality of the fitted models was evaluated using
the regression metrics defined in Appendix~\ref{app:evaluation_metrics}, namely the coefficient of
determination \(R^2\) and the Mean Absolute Error (MAE). 


\subsection{Runtime Estimation for Radius Neighbor Density}
\label{subsec:radius_runtime_estimation}
\hspace{\parindent}
To estimate the computation time of the \textit{Radius Neighbor Density} heuristic, an experimental
dataset was collected by measuring its runtime for different numbers of robots. The tested values
were \(N=\{1,5,10,20,50,100,150,200,300,400\}\), and, for each value of \(N\), the heuristic was
executed 50 times. The measured execution times were then used to fit different regression models.

The choice of candidate models was guided by Equation~\ref{eq:radius_neighbor_density_complexity}.
Accordingly, different regression models were tested using combinations of \(N\), \(N^2\), and
\(N\log N\). Table~\ref{tab:radius_runtime_models} summarizes the obtained results.

\begin{table}[H]
\centering
\caption{Regression models tested for estimating the computation time of the \textit{Radius Neighbor Density} heuristic.}
\label{tab:radius_runtime_models}
\small
\renewcommand{\arraystretch}{1.25}
\resizebox{\linewidth}{!}{%
\begin{tabular}{lccccc}
\hline
\textbf{Model} &
\(\boldsymbol{\theta^{\mathrm{rad}}_0}\) &
\(\boldsymbol{\theta^{\mathrm{rad}}_1}\) &
\(\boldsymbol{\theta^{\mathrm{rad}}_2}\) &
\(\mathbf{R^2}\) &
\textbf{MAE (ms)} \\
\hline
\(T_{\mathrm{radius}}(N)=\theta^{\mathrm{rad}}_0+\theta^{\mathrm{rad}}_1N\)
& -6.225741 & 0.155011944 & -- & 0.913134 & 6.518654 \\

\(T_{\mathrm{radius}}(N)=\theta^{\mathrm{rad}}_0+\theta^{\mathrm{rad}}_1N^2\)
& 0.487353 & 0.000339267 & -- & 0.990712 & 1.297389 \\

\(T_{\mathrm{radius}}(N)=\theta^{\mathrm{rad}}_0+\theta^{\mathrm{rad}}_1N+\theta^{\mathrm{rad}}_2N^2\)
& \textbf{0.413287} & \textbf{0.001414811} & \textbf{0.000336415} & \textbf{0.990718} & \textbf{1.282169} \\

\(T_{\mathrm{radius}}(N)=\theta^{\mathrm{rad}}_0+\theta^{\mathrm{rad}}_1N\log N+\theta^{\mathrm{rad}}_2N^2\)
& 0.467381 & 0.000089134 & 0.000338135 & 0.990713 & 1.293423 \\
\hline
\end{tabular}
}
\end{table}

As expected, the models containing the \(N^2\) term achieved the best results, clearly outperforming
the model based only on \(N\). This behavior is consistent with the complexity expression derived
in Equation~\ref{eq:radius_neighbor_density_complexity}, where the dominant term is associated
with the quadratic growth of the number of robot comparisons.

Among the models containing a quadratic term, the best result was obtained by combining \(N\)
and \(N^2\), leading to the following runtime model:

\begin{equation}
    T_{\mathrm{radius}}(N)
    =
    \theta^{\mathrm{rad}}_0
    +
    \theta^{\mathrm{rad}}_1 N
    +
    \theta^{\mathrm{rad}}_2 N^2 .
    \label{eq:radius_runtime_model}
\end{equation}

Although the complexity expression in
Equation~\ref{eq:radius_neighbor_density_complexity} also includes an \(N\log N\) component,
the difference between the \(N+N^2\) and \(N^2+N\log N\) models was almost negligible. This is
consistent with the fact that the dominant runtime behavior is captured by the \(N^2\) term, while
the remaining contribution can be sufficiently approximated by the linear term over the tested range
of robot numbers.

Figure~\ref{fig:lr_radius} shows the fitted curve obtained with the selected model, which captures
the non-linear increase in computation time as the number of robots grows, consistently with the
quadratic dominant term identified for this heuristic.

\begin{figure}[H]
    \centering
    \includegraphics[width=0.65\linewidth]{lr_radius.png}
    \caption{Regression model fitted to the computation time of the \textit{Radius Neighbor Density} heuristic.}
    \label{fig:lr_radius}
\end{figure}


\subsection{Runtime Estimation for \(K\)-hop Neighbor Density}
\label{subsec:khop_runtime_estimation}
\hspace{\parindent}
For the \textit{\(K\)-hop Neighbor Density} heuristic, the runtime estimation had to account for the
expansion process performed when computing the topological neighborhood of each robot. In this
case, the computation time is mainly affected by the total number of node expansions performed
during the \(K\)-hop searches.

The total number of expansions performed during one execution of this heuristic is mainly affected
by the average DoF of the roadmap vertices, the neighborhood depth \(k_{\text{hops}}\), and the
number of robots \(N\). Varying the roadmap degree would require creating new maps or modifying
existing ones only for this analysis, so this parameter was kept fixed. The dataset was therefore
generated by varying \(k_{\text{hops}}\) and \(N\), using
\(k_{\text{hops}}=\{1,2,3\}\) and \(N=\{10,25,50,100,150,200\}\). For each combination, the
heuristic was executed 20 times, and the corresponding computation times were recorded.
The tested regression models considered the measured total number of node expansions, denoted
by \(E\), together with terms depending on the number of robots. Table~\ref{tab:khop_runtime_models}
summarizes the obtained results.

\begin{table}[H]
\centering
\caption{Regression models tested for estimating the computation time of the \textit{\(K\)-hop Neighbor Density} heuristic.}
\label{tab:khop_runtime_models}
\small
\renewcommand{\arraystretch}{1.25}
\resizebox{\linewidth}{!}{%
\begin{tabular}{lccccc}
\hline
\textbf{Model} &
\(\boldsymbol{\theta^{\mathrm{khop}}_0}\) &
\(\boldsymbol{\theta^{\mathrm{khop}}_1}\) &
\(\boldsymbol{\theta^{\mathrm{khop}}_2}\) &
\(\mathbf{R^2}\) &
\textbf{MAE (ms)} \\
\hline
\(T_{\mathrm{khop}}(N,E)=\theta^{\mathrm{khop}}_0+\theta^{\mathrm{khop}}_1E\)
& 0.814326 & 0.002516826 & -- & 0.851778 & 1.176020 \\

\(T_{\mathrm{khop}}(N,E)=\theta^{\mathrm{khop}}_0+\theta^{\mathrm{khop}}_1N+\theta^{\mathrm{khop}}_2E\)
& 1.432990 & -0.013862 & 0.002977 & 0.880253 & 1.076947 \\

\(T_{\mathrm{khop}}(N,E)=\theta^{\mathrm{khop}}_0+\theta^{\mathrm{khop}}_1N\log N+\theta^{\mathrm{khop}}_2E\)
& \textbf{1.310290} & \textbf{-0.002688199} & \textbf{0.003005392} & \textbf{0.884116} & \textbf{1.058590} \\
\hline
\end{tabular}
}
\end{table}

The best-fitting model for the \textit{\(K\)-hop Neighbor Density} heuristic combines the
\(N\log N\) term with the total number of node expansions \(E\):
\begin{equation}
    T_{\mathrm{khop}}(N,E)
    =
    \theta^{\mathrm{khop}}_0
    +
    \theta^{\mathrm{khop}}_1 N\log N
    +
    \theta^{\mathrm{khop}}_2 E .
    \label{eq:khop_runtime_model}
\end{equation}

This result is consistent with Equation~\ref{eq:k_hop_neighbor_density_complexity}, as the expansion
term \(E\) captures the cost of computing the \(k\)-hop neighborhoods, while the \(N\log N\) term
models the remaining operations dependent on the number of robots, including the sorting step.
This explains why the model using both terms outperformed the alternative based only on \(E\).

Figure~\ref{fig:lr_khop} shows the fitted regression plane obtained with the selected model, using
a three-dimensional representation due to the dependence of the estimated computation time on two
input variables, \(N\log N\) and \(E\).

\begin{figure}[H]
    \centering
    \includegraphics[width=0.8\linewidth]{lr_khop.png}
    \caption{Regression model fitted to the computation time of the \textit{\(K\)-hop Neighbor Density} heuristic.}
    \label{fig:lr_khop}
\end{figure}



To use this model before executing the heuristic, the total number of node expansions \(E\) must be
estimated from the parameters of the current planning instance. In practice, the neighborhood
depth \(K\) is fixed for a given run, and high values of \(K\) are not typically useful in this context,
as they rapidly expand the considered neighborhood and may reduce the locality of the density
estimate.

For this reason, the expansion effort was approximated from the number of vertices reachable
within \(k_{\text{hops}}\) hops, previously defined in
Equation~\ref{eq:khop_reachable_vertices}. Assuming an average roadmap degree of \(DoF=4\),
this provides a practical estimate of the number of vertices explored around each robot. Although
the recursive implementation may expand some vertices more than once, this becomes unlikely for
the small values of \(k_{\text{hops}}\) considered here, making the approximation sufficient for runtime
prediction.

As the \(K\)-hop expansion is performed for each robot, the total expansion estimate used by the
runtime model is given by

\begin{equation}
    \hat{E}
    =
    N \cdot N_K
    =
    N \cdot \left(1 + 2k_{\text{hops}}(k_{\text{hops}}+1)\right),
    \quad DoF=4.
    \label{eq:khop_estimated_expansions}
\end{equation}

This estimated value \(\hat{E}\) is then provided to
Equation~\ref{eq:khop_runtime_model}, allowing the computation time of the
\textit{\(K\)-hop Neighbor Density} heuristic to be predicted before executing it.


\subsection{Estimated Runtime Comparison}
\label{subsec:estimated_runtime_comparison}
\hspace{\parindent}
After fitting both runtime models, their estimates were compared to select the faster heuristic for
each planning instance. Although the \textit{\(K\)-hop Neighbor Density} model also depends on
\(\hat{E}\), this term becomes a function of \(N\) once \(DoF\) and \(k_{\text{hops}}\) are fixed, as
shown in Equation~\ref{eq:khop_estimated_expansions}. Therefore, both models can be evaluated
as functions of \(N\) and compared directly.
Figures~\ref{fig:radius_khop_K1}, \ref{fig:radius_khop_K2},
\ref{fig:radius_khop_K3}, and \ref{fig:radius_khop_K4} show the estimated runtime comparison
for different values of \(k_{\text{hops}}\), assuming \(DoF=4\). The intersection between the two
curves indicates the approximate robot count at which both heuristics are expected to have the
same computation time.


\begin{figure}[H]
    \centering

    \begin{subfigure}[t]{0.48\linewidth}
        \centering
        \includegraphics[width=\linewidth]{K_1_comparison.png}
        \caption{\(k_{\text{hops}}=1\)}
        \label{fig:radius_khop_K1}
    \end{subfigure}
    \hfill
    \begin{subfigure}[t]{0.48\linewidth}
        \centering
        \includegraphics[width=\linewidth]{K_2_comparison.png}
        \caption{\(k_{\text{hops}}=2\)}
        \label{fig:radius_khop_K2}
    \end{subfigure}

    \vspace{0.4cm}

    \begin{subfigure}[t]{0.48\linewidth}
        \centering
        \includegraphics[width=\linewidth]{K_3_comparison.png}
        \caption{\(k_{\text{hops}}=3\)}
        \label{fig:radius_khop_K3}
    \end{subfigure}
    \hfill
    \begin{subfigure}[t]{0.48\linewidth}
        \centering
        \includegraphics[width=\linewidth]{K_4_comparison.png}
        \caption{\(k_{\text{hops}}=4\)}
        \label{fig:radius_khop_K4}
    \end{subfigure}

    \caption{Estimated computation time comparison between \textit{Radius Neighbor Density}
    and \textit{\(K\)-hop Neighbor Density}, considering different values of \(k_{\text{hops}}\)
    and \(\mathrm{DoF}=4\).}
    \label{fig:radius_khop_comparison}
\end{figure}



The intersections occur approximately at \(N=54.47\), \(N=99.94\), \(N=191.62\), and
\(N=326.65\) for \(k_{\text{hops}}=1,2,3,\) and \(4\), respectively. This increasing threshold is
consistent with the fact that larger \(k_{\text{hops}}\) values require more neighborhood
expansions, delaying the point at which the \textit{\(K\)-hop Neighbor Density} heuristic becomes more
efficient than the radius-based alternative.

Thus, for each fixed value of \(k_{\text{hops}}\), the corresponding intersection defines the
selection threshold used by the planner. Below this threshold, the \textit{Radius Neighbor Density}
heuristic is expected to be faster, while above it the \textit{\(K\)-hop Neighbor Density} heuristic becomes
preferable due to the quadratic growth of the radius-based model. Therefore, before launching the
worker threads, both fitted regression models are evaluated for the current planning instance, and
only the heuristic with the lowest estimated computation time is inserted into the heuristic pool.


\section{Summary}
\label{sec:parallel_summary}
\hspace{\parindent}
This chapter presented the parallel planning framework developed around TEA*. The proposed
framework builds on a previous multi-threaded TEA* prototype, where parallel execution was used
to explore multiple robot priority permutations~\cite{ribeiro2025parallel}. In this dissertation, that
idea was redesigned into a dedicated TEA* Parallel Planner class and extended with the
heuristic-based ordering methods developed in Chapter~\ref{chap:ch5}. Each worker thread runs an
independent TEA* instance, producing candidate plans that are later compared according to the
selected cost function.

OpenMP was adopted in the redesigned framework to launch the main worker threads, while
synchronization is limited to shared structures such as the heuristic pool. The mutable state of each
TEA* instance remains private to its worker, avoiding frequent locking during the most
computationally intensive parts of the planning process. This increases memory usage, since each
TEA* instance maintains its own mutable planning structures, including the 3D costmap used to
store space-time reservations. Nevertheless, this trade-off was considered acceptable because it
reduces contention between threads and helps preserve the benefit of evaluating multiple priority
orderings concurrently.

The chapter also detailed the heuristic dispatching strategy used to assign priority orderings to the
worker threads. The proposed mechanism relies on a shared priority queue organized by heuristic
groups, promoting a balanced use of the available heuristic families and reducing the probability of
evaluating redundant orderings in parallel. The \textit{Exploration Conflicts Pressure} heuristic was
also integrated as a fallback strategy when a worker exhausts the regular heuristic pool without
finding a valid plan.

The timeout and controlled termination mechanism was then described, showing how the planner
uses the available planning time without stopping immediately after the first valid solution is found.
Finally, a regression-based selection mechanism was introduced to handle redundancy between the
\textit{Radius Neighbor Density} and \textit{\(K\)-hop Neighbor Density} heuristics. Since these heuristics
often produce equivalent or very similar priority orderings, only the alternative with the lowest
estimated computation time is inserted into the heuristic pool.